\documentclass[a4paper,12pt,twoside,openright,titlepage]{article}

%Additional packages
\input{styles/packages}

% Listing style
\input{styles/lstset}

% Uncomment for production
% \syntaxonly

% Style
\pagestyle{headings}

% Turn on indexing
\makeindex[intoc]

% Define document
\author{D. Leeuw}
\title{Linux: Editors}
\date{\today\\
1.1.0\\
\vfill
\raggedright
\copyright\ 2020-2026 Dennis Leeuw\\
\input{styles/licentie-titlepage}}


\begin{document}
\selectlanguage{dutch}

\maketitle

%%%%%%%%%%%%%%%%%%%
%%% Introductie %%%
%%%%%%%%%%%%%%%%%%%

\section{Over dit Document}
\subsection{Leerdoelen}
\input{src/leerdoelen_cli_editors}
\subsection{Voorkennis}
\input{src/voorkennis_cli_editors}

%%%%%%%%%%%%%%%%%
%%% De inhoud %%%
%%%%%%%%%%%%%%%%%

% Requires: fhs
% Provides: vi, vim, nano, pico
\section{Het gebruik van een editor}
Op de commandline heb je geen menu's en vaak geen muis om door een applicatie te navigeren. Het maken van documenten is dan ook een stuk lastiger dan in een grafische interface. Toch zijn er oplossingen om op de commandline te werken met bestanden. Een tekstverwerker\index{tekstverwerker} op de commandline heet een editor\index{editor}. Met een editor kan je geen documenten opmaken, je kan er alleen platte tekstbestanden mee schrijven en bewerken. Er zijn in de loop der jaren verschillende editors bedacht. Een van de oudste voor Unix geschreven editors is \texttt{vi}.

Een van de grote namen achter Unix is Ken Thompson. De eerste drie commando's die hij schreef voor het jonge Unix systeem waren \texttt{as} (assembler), \texttt{ed}\index{ed} (editor) en \texttt{sh} (shell). Dennis M. Ritchie bracht verbeteringen aan op \texttt{ed} en vanaf 1969 tot 1976 bleef dit de editor op een Unix systeem. In 1976 kwamen Billy Joy en Chuck Haley met een nieuwe editor die \texttt{ex}\index{ex} werd genoemd. Voor \texttt{ex} schreef Billy Joy ook een soort interface om er makkelijker mee te kunnen werken en die wrapper om \texttt{ex} noemde hij \texttt{vi}\index{vi} (visual interface\index{visual interface}). Vanaf 1979 werd \texttt{ex} geintergreerd in \texttt{vi} en was er alleen nog \texttt{vi}. Later werd \texttt{vi} onderdeel van de Single Unix Specification en daarmee een editor die op bijna elk Unix systeem aanwezig is en dat is nogsteeds het geval. Op bijna alle beschikbare Unix systemen, van BSD tot Linux en Mac OS X is een vorm van \texttt{vi} aanwezig. Dat is dan ook het voordeel van het aanleren van het werken met \texttt{vi}, je kan je kennis op verschillende platformen gebruiken. Het nadeel van \texttt{vi} is dat het in het begin lastig is om ermee te leren werken.


\subsection{vi, pico, nano}
De eerste redelijk gebruiksvriendelijke editor op Unix was \texttt{vi}. De \texttt{vi} editor kent twee modi. De eerste modus is de \textquote{edit mode} en de tweede is de \textquote{command mode}. Standaard start \texttt{vi} op in de command mode waarin je commando's kunt geven om bestanden te laden of op te slaan en waarin je functies als knippen en plakken kan uitvoeren. De edit modus is die waarin je tekst invoert. Dit onderscheid maakt het voor beginnende gebruikers soms verwarrend.

Naast \texttt{vi} zijn er ook andere editors voor Unix-achtige systemen ontwikkeld. De meeste bekende zijn \texttt{pico}\index{pico} en \texttt{nano}\index{nano}. Pico was de oorspronkelijke editor. Nano is ontwikkled door het GNU-project en is een vervanging van pico omdat pico een licentie had die \textquote{problematisch} was. Dat probleem is inmiddels opgelost, maar nano biedt zoveel extra mogelijkheden dat velen de voorkeur geven aan nano.

Het grote voordeel van nano ten opzichte van vi is zijn gebruiksvriendelijke interface. Nano kent geen edit en command mode zoals vi. Nano gebruikt control codes om commando's te geven en is direct beschikbaar voor de invoer van tekst van de gebruiker.


\section{nano}
GNU nano is gebruiksvriendelijke en makkelijk aan te leren editor. GNU nano gebruikt \keys{CTRL} en de \keys{ALT} toets met een letter om commando's uit te voeren zoals opslaan van een document. De meest gebruikte toets-combinaties staan onder aan het scherm, alle toets-combinaties kan je terug vinden in de handleiding die je in \texttt{nano} kan opvragen met \keys{CTRL + G}.


\subsection{Knippen en plakken}
Om tekst te selecteren, verplaatst je de cursor naar het begin van de tekst en druk je op \keys{ALT + A}. Als je nu de cursor verplaatst met de pijltjes toetsen dan zal er tekst geselecteerd worden. Verplaats de cursor met de pijltoetsen naar het einde van de tekst die je wilt selecteren.

Met \keys{ALT + SHIFT + 6} kopieer je de tekst. Wil je de tekst knippen dan gebruik he \keys{CTRL + K}

Om de tekst te plakken, verplaats je de cursor naar de plek waar je de tekst wilt plaatsen en druk je op \keys{CTRL + u}.



\section{vim}
De naam \texttt{vim}\index{vim} staat voor Vi IMproved\index{Vi IMproved}. Of wel een verbeterde versie van \texttt{vi}. Bram Molenaar, een Nederlandse software ontwikkelaar, schreef sinds 1991 tot aan zijn dood in 2023 aan de code van \texttt{vim} en zijn verbeteringen zijn zo populair dat op alle Linux systemen alleen nog \texttt{vim} ge\"installeerd wordt. Je kan op een Linux systeem \texttt{vim} opstarten als \texttt{vi} waarmee je zoveel mogelijk de functionaliteiten krijgt als het oude \texttt{vi} en je kan \texttt{vim} opstarten als \texttt{vim} waarmee je alle nieuwe toevoegingen van Bram Molenaar en zijn medeontwikkelaars krijgt.

\texttt{vim} is niet een van de makkelijkste editors maar heeft als grote voordeel, zoals eerder gezegd, dat het beschikbaar is op elk willekeurig Unix systeem.


\subsection{vim opstarten}
\input{src/vimOpstarten}
\subsection{Text invoer}
\input{src/vimTextInvoer}
\subsection{Text delete}
\input{src/vimDelete}
\subsection{Knippen en plakken}
\input{src/vimKnippenPlakken}
\subsection{Bewegen door een tekstbestand}
\input{src/vimBewegen}


%%%%%%%%%%%%%%%%%%%%%
%%% Index and End %%%
%%%%%%%%%%%%%%%%%%%%%
\printindex
\end{document}

%%% Last line %%%
